\chapter{Introduction}

\section{Context and motivation}

Much of the API documentation is written 
in natural language and is available in formats designed 
to be read by people, not machines. 
This makes automation difficult, since the information 
is often scattered across HTML portals, PDF documents, 
manually exported Postman collections, or incomplete 
OpenAPI specifications that are no longer up to date with respect to the 
actual behavior of the service. In many cases, 
there isn’t even a formal specification.

Reconstructing this information manually is a slow, 
repetitive process prone to human error. Therefore, 
automating this task is essential, as is 
having a tool that facilitates the generation and 
maintenance of up-to-date specifications.

This project automates this reconstruction process 
by using a local language model capable of 
extracting the structure of an OpenAPI specification 
from unstructured text documentation. However, 
the central aspect of the design is to strike a 
balance between the system’s ability to extract information 
and its reliability, avoiding blind reliance 
on the responses generated by the language model. 
To this end, validation mechanisms are incorporated that 
allow for verifying the consistency and quality of the 
generated specifications before accepting them 
as the final result.


\section{General objective}

Build a system that functions as a tool capable 
of processing API documentation, regardless of 
its source format whether HTML, PDF, Markdown, 
Postman collections, or partial OpenAPI specifications and reconstruct 
a formal, verifiable model of the API in the form of a 
valid OpenAPI 3.x specification. Based on this 
specification, the system automatically generates three 
reusable artifacts, including a client SDK, 
a mock server, and a contract testing suite, 
thereby facilitating the integration, validation, 
and use of the API.


\section{The loop of verification}
Due to the phenomenon of hallucinations in language models, 
the process of transforming documentation into an OpenAPI 3.x 
specification can produce incorrect results, such as the creation 
of nonexistent endpoints, the erroneous assignment of data types, 
or the duplication of schemas. In this context, the project’s goal 
is to implement mechanisms that detect and mitigate these 
inconsistencies using structural validation techniques, 
rule-based linting, and dynamic probing when access to the actual 
API is available. The discrepancies identified during these stages 
are used as feedback to iteratively refine the generated 
specification, increasing its accuracy and consistency with 
the service’s actual behavior.


\section{Scope and Approach}
The project is conceived as a personal endeavor 
and is based on two preliminary design decisions that guide
its implementation and execution.

\subsection{Verification with graceful degradation}
During the project’s execution, an additional dynamic probing 
module is developed, whose function is to compare the generated 
OpenAPI specification with the API’s behavior during execution, 
provided that a sandbox environment is available for testing. 
This process allows for a more representative fidelity metric by 
evaluating the correspondence between the reconstructed 
specification and the service’s actual behavior.

In cases where a sandbox environment is not available, 
the verification process relies exclusively on structural 
validation and linting mechanisms, ensuring that the 
specification complies with the syntactic requirements and 
best practices defined for OpenAPI, although without validating 
its conformity with the actual implementation of the API.

\subsection{Target API with ground truth}
An API with an official, publicly available OpenAPI specification
such as those for Swagger Petstore, Stripe, or GitHub is 
selected as the evaluation target. During the reconstruction 
process, the system is provided only with the API documentation, 
while the official OpenAPI specification is deliberately hidden 
to prevent it from influencing the generated results.

Once the reconstruction is complete, the resulting specification 
is compared with the official specification, which serves as the 
ground truth. This comparison allows for the calculation of a 
fidelity metric that quantifies the degree of correspondence 
between the reconstructed specification and the original, 
providing an objective measure of the system’s accuracy.

\section{General architecture of the system}

\subsection{Stages of work}

The system is structured as a five-stage pipeline.
An intermediate model is exchanged between each stage, this model 
is implemented using a set of Pydantic models that act as a data 
contract between the system’s various components. In this way, 
each stage consumes the model generated by the previous stage 
and produces the model for the next stage, without needing to 
know the internal implementation details of the other components.

\begin{enumerate}[label=\arabic*.]
    \item \textbf{Ingestion and normalization} - This stage composes from the
    raw source and conversion to tagged fragments for each candidate endpoint
    \item \textbf{LLM extraction} - Each fragment is extracted, with a structured
    output, that in this project is an \texttt{EndpointSpec}.
    \item \textbf{Assembly} - The loose \texttt{EndpointSpec} merge
    in a document OpenAPI 3.x, deduplicating schemas and resolving references
    \item \textbf{Verification} - Structural validation, linting,
    and dynamic probing.
    \item \textbf{Generation} - From the verified spec it produces the SDK client,
    the mock server and the tests suite
\end{enumerate}

\subsection{The intermediate model and repository}

The repository structure reflects the stages of the pipeline
(\texttt{ingest/}, \texttt{extract/}, \texttt{assemble/}, \texttt{verify/},
\texttt{generate/}) plus a \texttt{models/} module that contains the 
shared contract among the different stages and processes of the pipeline.
The goal of this is to keep each module testable in isolation. 

\section{Technological Stack}

\begin{table}[h!]
\centering
\begin{tabular}{@{}ll@{}}
\toprule
\textbf{Stage} & \textbf{Tools} \\
\midrule
Ingest & \texttt{markitdown}, direct JSON parsers \\
Extraction & LLM with forced structured output, Chroma/Qdrant (RAG) \\
Assembly & Serialization Pydantic $\rightarrow$ OpenAPI 3.x \\
Verification & \texttt{openapi-spec-validator}, Spectral, Schemathesis \\
Generation & \texttt{openapi-generator-cli}, Prism, Schemathesis/Dredd \\
Orchestration & Custom state machine \\
\bottomrule
\end{tabular}
\caption{Technological stack per pipeline stage}
\end{table}

\section{Document structure}
This document is organized according to the project's development phases, 
with each chapter corresponding to a phase; it documents the phase's objective,
the design decisions made, the relevant code, and evidence that the phase deliverable
has been completed. This chapter covers exclusively \textbf{Phase 0 - Staging}; 
the subsequent phases (Ingest, Extraction, Assembly, Verification, and Generation) 
are documented in subsequent chapters as they are completed.

